\section{Prime Agent Architecture}
\label{sec:architecture}
We next describe the Prime Agent architecture in detail.
In particular, we discuss (1) how the system manages information and computation (\S\ref{sec:architecture-overview}), (2) how state is organized across model weights, active context, the REPL, and disk-backed storage (\S\ref{sec:information-hierarchy}), (3) how persistent REPLs and RLM calls support programmatic computation (\S\ref{sec:programmatic-computation}), (4) how recursive sessions communicate with agents and human operators (\S\ref{sec:recursive-orchestration}), (5) how Continual Harness turns trajectory evidence into reusable state (\S\ref{sec:continual-harness}), and (6) how long-horizon controls define continuation, termination, and evaluation accounting (\S\ref{sec:long-horizon-execution}).
\Cref{fig:architecture} provides an overview.

\begin{figure}[!t]
\centering
\includegraphics[width=\linewidth]{figures/prime_agent/architecture.pdf}
\caption{Prime Agent connects persistent root and subagent sessions to a daemon, Continual Harness, the Agents View, and the environment. Solid arrows carry execution and messages; dashed arrows carry persistent state.}
\label{fig:architecture}
\end{figure}

\subsection{Architecture overview}
\label{sec:architecture-overview}
Prime Agent separates information management from computation management.
Information management determines what state enters a model invocation and what survives compaction or restart.

Computation management maps model-selected actions to code, tools, and recursive subagent sessions.
Direct agent-to-agent communication connects related sessions, and direct human-agent interaction exposes individual nodes for inspection and intervention.

Models manipulate intermediate values with code.
Sessions retain history across compaction, detachment, and restart.
Subagents inherit the root's execution and communication primitives.
The runtime records model calls, tool use, messages, harness changes, and resource use.
The model controls decomposition, computation allocation, communication, and stopping.



\subsection{Information hierarchy and persistent state}
\label{sec:information-hierarchy}
A model invocation is parameterized by fixed weights and conditions on its active token context.
Prime Agent adds state outside that context. External state affects generation when the runtime injects it or an operation serializes a result into the context.
\Cref{fig:llm-hierarchy} organizes this state by visibility, access mechanism, and persistence.


\begin{figure}[H]
\centering
\includegraphics[width=\linewidth]{figures/prime_agent/state-hierarchy.pdf}
\caption{\textbf{Prime Agent state hierarchy.} The boundary between L1 and L2 separates token-visible model state from explicitly managed computation and retained state.}
\label{fig:llm-hierarchy}
\end{figure}

Each level changes through a different mechanism.
Fine-tuning updates L0, compaction rewrites L1, and refinement versions selected L3 entries.
We call the L2 mechanism \emph{agentic garbage collection}.
The model creates, retains, summarizes, or deletes REPL values and subagent sessions as the task changes.

Explicit operations move information between levels.
Python values and tool outputs in L2 enter generation when serialized into L1.
Compaction replaces a conversational prefix with a summary and retains the original events in L3 for REPL retrieval.
The runtime assembles selected Continual Harness entries into later supplemental prompts; other L3 artifacts enter context on retrieval.
L0 stays fixed.

The retained runtime state includes an append-only event history, selected kernel snapshots, the rooted session tree, context and compaction records, persistent message queues, and versioned Continual Harness state.
Branching or forking creates a new logical continuation without deleting the prior event sequence.
Recovery reconstructs the session under the same identity.
Non-serializable Python objects and external processes are recreated from saved artifacts or external services.

\subsection{Programmatic computation with RLMs}
\label{sec:programmatic-computation}
Each session owns a persistent IPython Read-Eval-Print Loop (REPL).
At test time, compute comprises model inference, Python execution, and tool calls; evaluations report tokens, time, and cost separately.
Installed tools are imported as Python modules for parsing, filtering, aggregation, and verification with ordinary code.
Intermediate values persist across turns and remain outside active context until selected.
This avoids repeatedly serializing large logs, task specifications, and structured evaluator output into the context.

Prime Agent implements the RLM abstraction with the asynchronous \rlmcall primitive~\citep{zhang2025recursivelanguagemodels}.
Calling \rlmcall creates and schedules a subagent session, then returns a stable handle before the subagent completes.
The subagent receives its own model context, IPython kernel, history, and workspace metadata.
The parent continues local computation while subagents run.
Results arrive later through direct agent-to-agent communication, and retained handles support follow-up after compaction or restart.
The model chooses between local code, tools, sequential delegation, and parallel subagents.
Prime Agent defines their execution semantics instead of a fixed workflow graph.
A complete orchestration example appears in \cref{app:orchestration}.

\subsection{Recursive orchestration and interaction}
\label{sec:recursive-orchestration}
The daemon owns live sessions independently of the client that created them.
Root and subagent sessions use the same lifecycle.
Sessions are running during a turn or tool operation, idle when loaded without an active turn, and inactive when unloaded but recoverable from persistent state.
Client detachment leaves the session running.
Stable session and parent identifiers preserve the recursive topology across these transitions.

Direct agent-to-agent communication uses asynchronous, daemon-mediated queues.
An agent addresses its parent, children, and siblings, and queued messages remain available when a recipient becomes active again.
Filesystem, network, and credential access follow the permissions of the runtime environment.
\Cref{fig:orchestration-lifecycle} summarizes the shared lifecycle and communication topology.

\begin{figure}[!t]
\centering
\includegraphics[width=\linewidth]{figures/prime_agent/orchestration-lifecycle.pdf}
\caption{Multi-agent orchestration lifecycle and direct agent-to-agent communication.}
\label{fig:orchestration-lifecycle}
\end{figure}

The Agents View exposes the persistent tree for direct human-agent interaction.
The interface lets a user inspect history, attach to a session, provide new input, or detach without interrupting execution.
The \code{agent-observe} interface provides bounded read-only status and recent-message previews; \code{agent-message} targets a named related session.
This allows for full interaction via the orchestrator.


\subsection{Continual Harness}
\label{sec:continual-harness}
Continual Harness exposes supplemental state for trajectory-time reads and writes~\citep{karten2026continualharness}.
Prompt notes store behavioral instructions, memories store facts, skills package executable procedures, and subagent specifications store reusable roles or divisions of labor.
Typed state separates rules, facts, programs, and coordination patterns.
Entries support create, read, update, and delete operations; local entries belong to one session, and explicitly requested global entries remain available to later sessions.

Refinement converts trajectory evidence into versioned state updates.
Agents request edits directly, or \code{/refine} runs a background model call over relevant events.
The runtime applies each edit at a turn boundary, records its trigger and intended effect, and assembles supplemental state for the next invocation.
Versions preserve provenance and enable rollback.
Refinement supplements the immutable base prompt without rewriting foundational policy.

\emph{Self-improvement} converts execution evidence into persistent harness state that changes later behavior while model weights remain fixed.
Useful computations become skills, repeated coordination patterns become subagent specifications, and corrected assumptions become memories or prompt notes.
The resulting trajectory record also provides training data for later model generations.

\subsection{Long-horizon execution and evaluation semantics}
\label{sec:long-horizon-execution}
Prime Agent exposes three long-horizon control mechanisms (\cref{fig:long-horizon-controls}).

\begin{figure}[!t]
\centering
\includegraphics[width=\linewidth]{figures/prime_agent/long-horizon-controls.pdf}
\caption{\textbf{Long-horizon control mechanisms.} Autonomous mode continues under an explicit budget until an end-condition test passes. Goals retain an objective across continuations until agentic completion. Heartbeats initiate cron or timed turns.}
\label{fig:long-horizon-controls}
\end{figure}

Autonomous mode continues model turns within an explicit budget and evaluates a task-specified end-condition test after each turn.
A failed test returns bounded output for another attempt; turn, token, and wall-clock limits stop execution.
A goal retains an objective across continuations and ends through agentic completion, when the agent marks the goal complete.
Heartbeats initiate turns on cron or timed schedules.

Evaluation configurations bind task and tool interfaces to model and provider settings, compaction and refinement policies, retry policy, completion gates, and resource limits.
Accounting aggregates the root and descendant sessions, so delegation remains visible in test-time cost.
Event history links model and tool calls, messages, interventions, retries, verifier outcomes, and harness edits to that configuration.
Standardized persistence, recovery, termination, and accounting separate harness failures from model failures while preserving model control over decomposition.
